\scanpage{229}
\mspage{117}
\deleted{L'opération de $Z$}

Considérons
\[
 S^{\unsure{\leftrightharpoons|\natural|\leftrightarrows|\flat}}=S/Z_S,\qquad S^{\vee\leftrightharpoons}=S^\vee/Z_{S^\vee}
 \tag*{(34)}
\]
alors les $\varphi\in\pi$ définissent par passage aux quotients
une $=$ application bijective
\[
 \varphi^\leftrightharpoons:S^\leftrightharpoons\xrightarrow{\sim}S^{\vee\leftrightharpoons}
 \tag*{(35)}
\]
par quoi on identifie $S^\leftrightharpoons,S^{\vee\leftrightharpoons}$ à un $=$ ensemble $B$.
\deleted{On pose de plus} Les orbites \inserted{\unsure{$Z$, ou de|$Z$-orbites}} de $Z_S$ dans $S$ s'appellent \uline{rayons}, \deleted{de \unsure{$B$},}
celles dans $S^\vee$ s'appellent corayons -- donc la donnée
de $\pi$ nous définit une bijection entre \inserted{ens. des} rayons et
\inserted{ens. des} corayons. On posera
\[
 R^\leftrightharpoons=R/Z \tag*{(36)}
\]
et \deleted{l'application $\subset$} \inserted{l'inclusion} $R\subset S\times S^\vee$ définit par passage
aux quotients
\[
 R^\leftrightharpoons\longrightarrow S^\leftrightharpoons\times S^{\vee\leftrightharpoons}\tag*{(37)}
\]
qui n'est plus injective en général, \deleted{En fait,} \inserted{mais qui l'est souvent.}
\deleted{L'opération de $Z$ sur $R$ \unsure{provient|procède|provenant} d'une opération de
$Z\times Z$ sur $S\times S^\vee$}
\deleted{sur} Supposons pour simplifier que $Z$ opère librement
sur $S$ \deleted{et} $S^\vee$ \inserted{(sur $S^\vee$ et sur la \unsure{relation})}
(une hypoth. implique l'autre d'ailleurs
quand on en a une en plus un $\pi$\ldots).
% slanted
\margin{hypothèse vérifiée dans le cas biséparé ``homogène'' \deleted{\unsure{qu'on suppose} \illegible{2}}}
Alors les rayons et les corayons sont des torseurs sous $Z$.

Soient \deleted{$\unsure{D\in S}$,}
\[
 \underset{\substack{\rotatebox{90}{$\in$}\\S^\leftrightharpoons}}D\subset S,\quad
 \underset{\substack{\rotatebox{90}{$\in$}\\S^{\vee\leftrightharpoons}}}{D^\vee}\subset S^\vee
\]
\deleted{deux} rayon et corayon ``liés'' i.e. \deleted{formant} correspondant
à un $(s^\leftrightharpoons,s^{\vee\leftrightharpoons})\in S^\leftrightharpoons\times S^{\vee\leftrightharpoons}$ dans l'image de $R^\leftrightharpoons$,
i.e. tels que
\[
 \exists s_0\in D,\ s_0^\vee\in D^\vee\quad\text{tels que }(s_0,s_0^\vee)\in R
 \qquad(\text{NB }D=Zs_0,\ D^\vee=Zs_0^\vee) \tag*{(38)}
\]
Alors la \deleted{partie} \inserted{fibre} de $R^\leftrightharpoons$ \deleted{sur} \inserted{en} $(s^{\leftrightharpoons},s^{\vee\leftrightharpoons})$ est le quotient
% archive sheet number: 228

